% 原书第 387 页（PDF 第 400 页）
\BookPage{387}{400}
\BookHeading{清代乾隆时期扬州人的引领时尚}{清代乾隆时期扬州人的引领时尚\\[3pt]—建设文化教育休憩城的历史启示}{冯尔康*}
十五年前，我为《扬州研究一江都陈轶群先生百年冥诞纪念论文集》作序，说到「扬州似乎常和『时髦』联系在一起，被作为繁华和奢靡的代名词」。又说「商品经济刺激人追求时髦，扬州的兴衰史似可证明这一点，这或许也是给人们的历史启示吧」。这是提出扬州城市的某种特点问题，未能探究，借这次「盛清社会与扬州」学术研讨会的机会，想谈谈这一历史启示，也是一定程度上了却一椿宿愿。

扬州在历史上的繁华，有三次高潮，一为隋炀帝下扬州，时间极其短暂；二为唐代，有「扬一益二」的美誉，作为中央政府财政（盐业度支）重镇和内河国际港口，风光百年；三为清朝前期，历经康熙后期、雍正乾隆二朝和嘉庆初年，而以乾隆年间为鼎盛，其时是以盐业产销为主体的商品经济发展，是运河、长江交通枢纽，贸易集散地，文化繁荣。后两次高潮的共同特点是经济与文化同步发展，令后人能够领略其时人们高度发达的社会文化生活行为和享受状况。

本文所说的扬州人，主要是指清代扬州府城郭的居住者，由于扬州府辖有泰州、高邮、仪征、兴化等州县，这些地方的居民生活亦兼有涉及。所说的时尚，是指文化生活、社会生活（衣食住行交游）行为，并形成风习；所谓引领，不仅是本地人如此生活，还能影响他处，将俗尚移植远方，成为他处人学习的目标和行为。

\BookSubheading{一、扬州最繁荣的时代在乾隆时期}

我想通过明清时期各界人士及近当代学者对扬州的评论，来认识扬州从经济发展演变为经济、文化双发达的城市，而在乾隆时期处于鼎盛状态。

\vfill
\begin{BookNotes}
\BookNote{*}{作者任职于南开大学中国社会史研究中心学术委员会。}
\BookNote{1.}{《扬州研究一江都陈轶群先生百年冥诞纪念论文集》，台北，联经出版公司，1996年。}
\end{BookNotes}

% 原书第 388 页（PDF 第 401 页）
\BookPage{388}{401}
明代嘉万间，文学家和史学家王世贞，以历史的、文化的眼光看待扬州，说唐代「广陵之富甲于诸藩镇，词人骚客又多为歌咏以益之，而古迹名胜灿然备矣」是说唐代扬州兼具经济、文化两种繁华，到了他生活的年代，经济复盛，「然而一盐客薮耳，其于大雅未复也」2。即经济发展之时，文化未能同步兴盛，留有遗憾。

到了清代，顺康之际在扬州任职推官的王士禛云:「近日地气自江南至江北，而扬州为极盛。」是扬州在明清之际战争疮痍之后，处于迅速恢复之中。及至乾隆帝六下江南，在扬州驻跸游览，其〈自高旻寺行宫再游平山堂即景六首〉之二写道，「从来富庶说广陵，满城丝管映街灯。康风拟令崇浮约，谋食贫人虑失屏」4。二十七年（1762）的〈平山堂绝句八首〉之二又云，「楼台丝管广陵擅」5。歌咏的是扬州一派升平繁华气象，尤以楼台丝管为说词，即不仅是经济的富庶，更是生活艺术化，表明人们生活中富有文化内涵。雍正间两淮盐政高斌、盐运使尹会一皆说，「广陵名郡，人文渊薮」。高斌还讲，扬州「人物秀杰、掇巍科，登显甲者，如云而起」7。直接说出扬州人才辈出、人文荟萃、文化发达的事实，当然他们深知这是经济发展促成的。两淮盐政阿克当阿说，「维扬，东南大都会」，「其人文之盛，尤史不绝书」8。仪征人李斗在乾隆二十九年（1764）至六十年（1795）的三十余年间，搜集扬州人活动的素材，于乾隆末年写出《扬州画舫录》，以亲身经历、耳闻目睹之事，反映扬州人生活的繁华景象。阮元于嘉庆二年（1797）为该书作序，谓乾隆帝「翠华六巡，恩泽稠叠，士日以文，民日以富」，所以〈扬州画舫录序〉系「目睹升平」之实录。他明确「士文」、「民富」是繁荣标志，扬州兼而有之。袁枚在〈扬州画舫录序〉中说：「扬州一郡，又为风尙华离之所。」自乾隆帝十六年（1751）南巡，官吏商民营建，扬州面貌较康雍时期大改观。袁枚居住南

\vfill
\begin{BookNotes}
\BookNote{2.}{王世贞，〈张孟奇〈广陵怀古诗〉序〉，嘉庆《江都县志》卷9〈艺文〉。}
\BookNote{3.}{王士禛，《香祖笔记》卷7，清刊本第15页b。}
\BookNote{4.}{见赵之璧，《平山堂图志》，第3页，《扬州地方文献丛刊》，广陵书社，2003年。}
\BookNote{5.}{《平山堂图志》，第4页。}
\BookNote{6.}{嘉庆《两淮盐法志》卷53〈杂纪二·书院〉。}
\BookNote{7.}{道光《重修仪征县志》卷16〈学校〉，光绪版。}
\BookNote{8.}{嘉庆《扬州府志·序》。}
\BookNote{9.}{李斗，《扬州画舫录》，中华书局，1980年。}
\end{BookNotes}

% 原书第 389 页（PDF 第 402 页）
\BookPage{389}{402}
\noindent 京，经常往来扬州，堂妹袁棠即嫁给扬州人，对扬州非常熟悉，他说此地是「风尙华离之所」，启发我们认识乾隆时代扬州的引领时尚。

繁华既然在于经济发展的民富，文化发达的士文。「士文」是什么意思？王世贞的说法是有士人名流活动形成的古迹名胜，即以名胜古迹为文化的标志。康乾时期扬州人、编修程梦星同意此说，并谓「扬州名胜之地」10。嘉庆十三年（1808）周长泰为顾銮的《广陵览古》作序，也说「广陵为古邗沟，名胜甲天下。竹西歌吹，至今传焉」11。都以扬州富有名胜而自豪。何谓名胜？扬州原来有隋炀帝冢、袁娘墓，亡国之君遗存不足称道，不可能成为名胜古迹的主要方面，文坛领袖王士禛在扬州的逸闻趣事，与同好的游览平山堂，留下诗文，「小秦淮」即出自其说，为扬州人津津乐道，给名胜增添色彩。康熙帝、乾隆帝南巡的临幸扬州，是最重大的盛事、幸事，扬州的官商为接驾而大肆营建，令市容大大改观。两淮转运使德庆在〈（嘉庆）扬州府志序〉谓该地，「况淮南大都，为两圣豫游之地乎！」皇帝去过的寺庙、林园身价倍增，平山堂、大明寺之类的不必说了，像乾隆朝编辑《四库全书》，扬州绅商马曰琯捐助私家藏书最多，乾隆帝嘉奖江浙人士的贡献，特在扬州建立文汇阁、杭州设文澜阁、镇江建文宗阁，赐书《古今图书集成》，及四库全书馆藏书的抄本，允许士子到阁借阅，文汇阁之名亦为钦赐。在扬州任安定书院院长的国子监祭酒吴锡麒创作〈广陵赋〉，特意写道：「独不见圣天子之惠我儒生乎？度恢宏而建阁，被『文汇』之嘉名。12正谊书院，明武宗正德间改祀董仲舒，康熙帝赐「正谊明道」额，遂名董子祠13。趣园，系盐商黄履暹别业，乾隆帝于二十七年（1762）临幸赐名，黄亦因接驾有功，恩赐奉宸苑卿4。这一个个寺观堂馆，有着帝王临幸，官员、文士流连的嘉话，更有歌咏图籍的记载。王世贞们所说的文化城市的内涵，或者说标准，可以分析出三点：外相的景观，名胜古迹，山水楼台；内在的是这些景观的人文含意，富有名人活动遗迹和遗文、逸闻，即有有趣的、有社会传播价值的故事，是可以传承的，是物质文化遗产与非物质文化遗

\vfill
\begin{BookNotes}
\BookNote{10.}{《平山堂图志》，第71页。}
\BookNote{11.}{见顾銮，《广陵览古》，《扬州地方文献丛刊》，广陵书社，2003年。}
\BookNote{12.}{同治《续纂扬州府志》卷23艺文志上〉。}
\BookNote{13.}{《扬州画舫录》，第207页。}
\BookNote{14.}{《平山堂图志》，第17页。}
\end{BookNotes}

% 原书第 390 页（PDF 第 403 页）
\BookPage{390}{403}
\noindent 产的综合体；三是歌舞艺术的行为与载体，这就是乾隆帝歌咏的「满城丝管映街灯」，兴化人郑板桥书写的：「画舫乘春破晓烟，满城丝管拂榆钱。」「扬州自古风流地，惟有当官不自怡。」15

明朝人谓扬州富而不文，满是铜臭气，予以蔑视；清朝人推崇扬州为人文渊薮的文化城市。是则扬州经过清朝前期、中期文化建设和文化积累，成为文化之乡，清人的见解是如实反映扬州实际状况的。

前述乾嘉道时期人们认为乾隆时代扬州是商业和文化重镇，近当代的学人认可这种说法。二十世纪上半叶，陈含光为《扬州丛刻》出版者陈恒和写墓志铭，内有「扬之地以文化称」之句16。清代扬州史的重要研究者王振忠在九十年代认为，乾嘉时期，「扬州成了东南首屈一指的文化中心城市。东南文化的精华都在此汇集、提炼和升华，并形成独具特色的文化内涵，进而向全国各地辐射、播布」17。杜瑜弟认为扬州在雍干时期「逐渐成为当时东南地区文化中心」，至近代才因淮盐转输和漕运的交通枢纽的优势丧失，发展不及他处18。卞孝萱在《郑板桥全集·前言》中云，「康雍干时期的扬州，是东南沿海地区一大都会。经济的繁荣，促成了文化艺术的昌盛。19

嘉庆后期，扬州的辉煌已逐渐褪去颜色。时至嘉庆十三年(1808)，目光敏锐的人已经感到扬州不如从前，于是说「扬城昔亦繁华耳」，「而今旧雨重相过，东道依然地主贫」20。旧友相逢，东道主的出手不似从前阔绰，显得寒酸了。无锡人钱咏多次盘桓扬州，于乾隆五十二年（1787）初次游览平山堂，看到的是「自天宁门外起直到淮南第一观，楼台掩映朱碧鲜新」，是令人赞叹的欣欣向荣景象；三十余年后再至，「几成瓦砾场，非复旧时光景矣」，不无伤感地说，「《（扬州）画

\vfill
\begin{BookNotes}
\BookNote{15.}{《郑板桥全集》，第30、70页，卞孝萱编，齐鲁书社，1985年。}
\BookNote{16.}{《扬州丛刻》于1930年前后印行。}
\BookNote{17.}{王振忠，〈明清徽商与扬州城市文化的特征和地位〉，载《扬州研究》，第502页。}
\BookNote{18.}{杜瑜弟，〈扬州历史地理综论〉，《扬州研究》第56-61页。}
\BookNote{19.}{卞孝萱，《郑板桥全集·前言》，见《郑板桥全集》。}
\BookNote{20.}{林苏门，《邗江三百吟》，第108页，《扬州地方文献丛刊》，广陵书社，2003年。}
\end{BookNotes}

% 原书第 391 页（PDF 第 404 页）
\BookPage{391}{404}
\noindent 舫录》中人半死，倚虹园外柳如烟」21。著名的马氏藏书、程氏藏砚的盛事，均成历史而不复存在。扬州府首县之一的甘泉县的陈家集（今属仪征），夙有小扬州之称，当地人林溥回忆：雍正乾隆间，镇中号全盛，收藏书画名迹极富，名流多往来于此；道光以来，富家式微，物力亦艰，灯市阑珊，迥非昔日矣22。然而瘦死骆驼比马大，嘉道以降的扬州虽与乾隆鼎盛之时无法相比，但市面并非一片萧条，仍有繁华余辉。

综观清朝人和当代人对扬州城市史的研究，我们可以认为乾隆时代扬州完全改变了明代纯粹商贾文化的粗俗面貌，在经济富庶同时，人文昌盛，创造新的名胜，赋予古迹新的文化内容，在生活享受中讲求艺术品味，令扬州以经济文化繁华城市的面貌出现于社会，扬州人的生活方式兼具文明与奢华两种成分，并能辐射四方。

\BookSubheading{二、扬州人引领时尚潮流的具体内容}

从文化方面看扬州的繁荣。然而事事物物皆有其文化内涵，我这里所讲的文化，排除商业文化的经营之道、官文化的做官之道，物质文化的生产与技术间或有所涉及，重点是在学术教育、文学艺术、信仰文化、消费文化、休闲文化。所说的消费文化，主要是指衣食住行交游的消费，还有正当的和狭邪的娱乐。具体内容来讲，有生活空间的住宅园林花圃、家庭用具摆设和工艺品；食物生存的餐饮，与此相关的人际交往的餐馆茶馆；美容健身的衣饰、沐浴、按摩、理发；口耳身心并赏的戏剧曲艺杂技；兴办学校与教育成效；狭义文化的经学、文学、史学研究，图书的刻印与收藏；精神寄托的宗教信仰、各种神灵信仰；社会救济及其观念；狭邪的娼妓、赌博；等等。是俗语的人生口福、眼福、耳福和心灵之福的一些具体事象。我试图将这些庞杂的内容加以归纳，无奈学术功力不济及研究的不足而不能做到。只好以下述九种类别，叙述扬州兴盛时代人们的生活及习尙，以及它的向外传播，至于如何认识它的价值，留待第三节交待。

\vfill
\begin{BookNotes}
\BookNote{21.}{《履园丛话》卷20<平山堂〉，第533页，中华书局，1979年。}
\BookNote{22.}{林溥，《扬州西山小志》，第22页，《扬州地方文献丛刊》，广陵书社，2003年。}
\end{BookNotes}

% 原书第 392 页（PDF 第 405 页）
\BookPage{392}{405}
\BookSubheading{（一）园林花卉甲天下}

这里关注的是兴盛时代的扬州人的生活空间情形：居室及其艺术，诸如住宅园林花圃，居家用物、摆设及工艺品。

乾隆时期扬州园林之胜，为天下之最。钱咏在其著作《履园丛话》说，「造物之工，当以扬州为第一」，素负盛名的苏杭不及。又说国初以张南垣为园林建筑高手，近时杰出者是常州人戈玉良，仪征巴氏朴园、江宁五松园皆其手笔，道光三年（1823）往游朴园，认为较苏州狮子园有趣，实为淮南第一名园23。吴锡麒在前述〈广陵赋〉说，扬州「园亭之胜，甲乎四方」。园林之盛，为扬州的特色。李斗说「郡城以园盛，康熙间有八家花园」，内有筿园、贺园。程梦星的筿园，营建于康熙末年，园内外百数十亩，建设有多处亭榭花圃，有荷花池和水榭的今有堂，梅园及亭的修到亭，种植芙蓉养蓄水鸟的月牙形池，竹丛中可眺可咏的来雨阁，三十株老桂树的桂坪，花药古松的馆松庵，别有竹畦，枯死，马曰琯赠竹，方士庶为绘赠竹图，程梦星遂名此园为筿园24。程氏「于艺事无所不能，尤工书画弹琴，肆情吟咏」。每当园花报放之时，辄与友朋同赏，作诗文之会，「以是推为一时风雅之宗」，受他延揽的有韦谦恒，探花出身，程氏「于家中构玉山心室，延之校书」。他的常客之一陈撰，中乾隆博学宏词科。盛唐，布衣，工书，「馆于筿园最久」。张铨，诸生，喜游山水，足迹遍天下，精于鉴别古人书画，主程氏，绘扬州二十四景及金、焦二山图画。程氏负时望，「江淮冠盖之冲，往来投赠殆无虚日」25。招徕之众，可以想见。贺园，贺君召建成十二景，接待四方宾客游览，两淮盐政准泰为其春玉堂题额「襟怀顿爽」，尚书、书法大家张照则题联：「万树琪花千圃药，一庄修竹半床书。」乾隆十一年（1746），贺氏大会宾客，蓄墨数升，供客人题咏，成诗千首。他汇集游人题咏的诗词，成《东园题咏》一书26。筿园、贺园等私家园林，对官绅友朋开放，形成逸闻趣事的嘉话，于是造物之工与人文雅事合成名

\vfill
\begin{BookNotes}
\BookNote{23.}{《履园丛话》卷10〈营造〉、卷12〈堆假山〉、卷20〈朴园〉，第326页、第330页、第534页。}
\BookNote{24.}{《扬州画舫录》，第22页、第343页。}
\BookNote{25.}{嘉庆《江都县续志》卷6〈人物〉，台北成文出版公司《中国方志丛书》本，华中地方第394号第1册第197页。}
\BookNote{26.}{《扬州画舫录》，第317-324页。}
\end{BookNotes}

% 原书第 393 页（PDF 第 406 页）
\BookPage{393}{406}
\noindent 胜，为其主人，也为地方增色。

园林宅院有花圃，而扬州别有花卉业，能与京城争胜，可知其不凡的地位。郑板桥的诗句「千家养女先教曲，十里栽花算种田」27，从种花之盛及其销路，道出花卉业的实况。扬州人种花是专业经营，并有专业行销者，消费者既有广大人群，又有专门行业的从业者。种花，不是一般的园圃业，而是专植花卉。傍花村、堡城、小茅山、雷塘等处均有集中种花者28。傍花村数十家，皆种菊为业，花时塡街绕陌，「秋间赏者如堵」29。徐谦芳述及扬州人的职业:四野之民业农，四郊之民业圃，沿江各洲之民业商，傍海之民业鱼盐，近郭之人以种花为业者亦多，莳芍药与菊，多佳种，且流布江南各处，几与北京丰台争胜30。反映花业盛行。莳花者中的宗元鼎，「手艺草花数十种，晨担花，向红桥坐卖，得钱沽酒，市人笑为花颠」。自着《卖花老人传略》31，是颇有个性的花匠。另一个来自苏州的移民，「手种奇花供客赏，三春一直到三秋」32。种花有自卖的，如宗元鼎，也有专业贩花发售的，与花农自卖者共同营造出繁盛的花市。他们每旦入城聚卖于市，如开明桥。花市有特色，如虹桥花市，二月梅花、桃花，夏季为牡丹、芍药、荷花，秋天为桂花、芙蓉33。花市之外，更有小贩沿街叫卖，郑板桥诗句「卖花声里雨蒙蒙」，「小楼帘卷卖花声」34，反映的就是此种情形。用花者甚多，所谓「扬人无贵贱皆戴花」35。「郭外饱看花，道观天桃争秀色，宝城秋菊斗奇葩，佳日赏芳华。」36记录扬州人的欣赏花卉情景，非只是社会上层人士的行为，普通居民亦参与其事。江都南门外花院莳茉莉、珠兰、白兰香橡花之类，专为贩户探买，制成花表等品，或穿花茶、供碟、花篮，制为三星桌围等物，以备礼品，转售平康乐户

\vfill
\begin{BookNotes}
\BookNote{27.}{《郑板桥全集》，第37页。}
\BookNote{28.}{《扬州画舫录》，第23页、第80页。}
\BookNote{29.}{林苏门，《邗江三百吟》，第13页。}
\BookNote{30.}{徐谦芳，《扬州风土记略》，第47页，《扬州地方文献丛刊》本，广陵书社，2003年。}
\BookNote{31.}{《平山堂图志》，第83页。}
\BookNote{32.}{董伟业，《扬州竹枝词》，第8页。}
\BookNote{33.}{《扬州画舫录》，第251页。}
\BookNote{34.}{《郑板桥全集》，第130页、第140页。}
\BookNote{35.}{《扬州画舫录》，第80页。}
\BookNote{36.}{惺庵居士，《望江南百调》，第1040页。}
\end{BookNotes}

% 原书第 394 页（PDF 第 407 页）
\BookPage{394}{407}
\noindent 及闺阁媛秀，四季无间37。妓院是探买的大户，闺阁媛秀亦喜爱有加。女性晚妆用花，如《望江南百调》所咏：「浴罢晚梳妆，雪白罗襦禅翼薄，霞青纱褂麝兰芳，花插夜来香。」38文士鉴赏鲜花，寄寓情怀，达到物我交流的境地。因为花业盛行，「花朝」也就过得有声有色，届时举行百花会，在张秀才家进行，四乡名花毕集献艺3。

家庭摆设、用物，扬州人讲究使用工艺品，漆器、玉器、剪纸、盆景中外驰名。据惺庵居士认知，扬州漆器「最精纯」，卢姓漆工用铜胎制作仿古器物，朱姓善于翻新花样4°。洋漆壁瓶，插花最佳，扬州驰名41。夏姓漆工善制古漆器，有剔红、塡漆两种，以金银铁木为胎，朱漆三十六次，镂以细锦，漆合有蔗段、蒸饼、河西三撞两撞诸式，漆盘有方圆、八角、牡丹花瓣诸式，漆器匣有长方两三撞诸式，统称为雕漆器42。面塑，袁姓业者最佳，捏人像能传神:「妙肖传神真面目，装潢余事小琴樽，绝技擅专门。」43牙雕，代不乏人，故清季有于啸轩者，技艺最为嘉善，方寸之间，能刻五千言以上；又能模刻名人山水、人物于牙屏之上，欧美万金争购44。荷兰女王加冕，湖广总督端方送刻牙一方，即于啸轩为之，正面是女王像，背面刻《金刚经》全文45。

\BookSubheading{（二）维扬菜系与发达的餐饮业}

扬州茶馆被称作天下第一，饭菜制作形成维扬菜系，予人口福，餐馆、茶馆不仅是适应人的生理需要，同时是人们社交的公共场所，具有多种功能。

\vfill
\begin{BookNotes}
\BookNote{37.}{《扬州风土记略》，第51页。}
\BookNote{38.}{《望江南百调》，第1049页。}
\BookNote{39.}{《扬州画舫录》，第80页。}
\BookNote{40.}{《望江南百调》，第1050页。}
\BookNote{41.}{林苏门，《邗江三百吟》，第39页。}
\BookNote{42.}{《扬州画舫录》，第198页。}
\BookNote{43.}{《望江南百调》，第1048页。}
\BookNote{44.}{《扬州风土记略》，第91页。}
\BookNote{45.}{杜召棠，《惜余春轶事》，第33页，《扬州地方文献丛刊》本，广陵书社，2003年。反映的虽系晚清之事，然而亦可以想像牙雕业早先的雄厚基础。}
\end{BookNotes}

% 原书第 395 页（PDF 第 408 页）
\BookPage{395}{408}
扬州食材丰富，郑板桥诗云:「百六十里荷花田，几千万家鱼鸭边。…湖上买鱼鱼最美，煮鱼便是湖中水。…一塘蒲过一塘莲，荇叶菱丝满稻田。最是江南秋八月，鸡头米赛珍珠圆。」「昨夜村灯鱼藕市，青帘醇酒见人情。」46原料及制作成的美食诱人，溢于纸上。惺庵居士的《望江南百调》记叙扬州名菜名点有烤鸭(「肥烤鸭皮包饼夹，浓烧猪肉醮馒头，口福几生修」）、三鲜面、煮干丝（「加料干丝堆细缕」）、水晶肴、蜜饯、酥糖（「蜜饯溅牙桃杏脯，酥糖到口桂兰香，风味最难忘」）等。酒肆甚多，在虹桥，康熙间有野园、冶春社、七贤居、且停车，为游人小酌之地，韩园则为聚饮之所。传统社会的厨师，分家厨、受人临时雇用的外厨、餐馆的行厨三种，据李斗的纪录，烹饪之技，家厨最胜，各家有绝活，如吴一山炒豆腐，田雁门走炸鸡，江郑堂十样猪头，施胖子梨丝炒肉，张四回子全羊47。维扬名菜「三头宴」的闷猪头，原来是学者江藩家的拿手菜，由来久远了。扬州点心制法极精，汤包、油糕尤擅名一时48。袁枚在《随园食单》中记录了许多扬州菜，如剥壳蒸蟹、程立方豆腐、定慧庵冬瓜、两淮盐运使卢见曾家的运司糕、洪府粽子，仪征萧美人点心，人参笋的做法是细切笋如人参形，微加蜜水，因「扬州人重之，故价颇贵」49。袁枚是美食家，他的品评多为定论。扬州菜在讲求色香味之外，别加装饰物，夏令厨师雕诸瓜为灯，玲珑透露，精美绝伦°。扬州人讲究食鲋鱼，四五月间产在扬子江，上市之时，乾嘉之际一尾价值三五千钱51。崇尚食卜鱼，故其价昂。鳉鱼的吃法，《调鼎集》罗列了十几种，有煮、蒸、红煎、淡煎之别，还可做成鳉鱼圆、鳉鱼豆腐、醉鳉鱼、糟鳄鱼、鳄鱼烩素面、鳄鱼面、鲋鱼羹，以及煨三鱼2。扬州盐商习于浮华，精于肴馔，造成扬州宴席各地驰名，及至扬州已衰落，而饮馔之盛尤甲于江苏53。

\vfill
\begin{BookNotes}
\BookNote{46.}{《郑板桥全集》，第63页、第127页。}
\BookNote{47.}{《扬州画舫录》，第26页、第253页。}
\BookNote{48.}{《扬州风土记略》，第48页、第91页。}
\BookNote{49.}{袁枚，《随园食单》，中国商业出版社，1984年。}
\BookNote{50.}{《扬州风土记略》，第91页。}
\BookNote{51.}{《邗江三百吟》，第126页。}
\BookNote{52.}{《调鼎集》，第365-369页，中国商业出版社，1986年。}
\BookNote{53.}{《扬州风土记略》，第86页。}
\end{BookNotes}

% 原书第 396 页（PDF 第 409 页）
\BookPage{396}{409}
李斗说:「吾乡茶肆，甲于天下。」54他是「三至粤西，七游闽浙，一往楚豫，两上京师」55的游历家，见多识广，不会作无根之谈，扬州茶馆应系甲天下。扬州城乡皆有茶馆，所谓「扬城四面多茶寮，由朝至暮，辄高朋满座，抵掌谈天，故北方人谓扬人为『渴相如』（渴死鬼之意）」56。乡镇亦有茶馆，甘泉「陈家集人每早必上茶馆，亦居然扬例也」57。惺庵居士说:「高会谷林堂，试茗有泉烹白雪。」反映扬州人的茶趣，人们举行茶会，品尝天下第五泉烹制的茶水。扬州人有到茶馆饮早茶的习惯，嘉庆以前，人们清晨相约，至好茶馆、好面馆，或叙谈，或议事58。品茗，为的是叙事闲谈，或者是商谈业务，且有点心供客人果腹，如林妪的茶肆，供应酥儿烧饼，见称于市9。茶馆有雅致的，林妪的茶馆，二门内有三间开的厅，题曰「秋阴书屋」，厅后住房十数间，二层，一层为客座，二层为卧室，因其近水、依城，「游人无不适意」。北门桥的双虹楼，铺面五开间，东壁开窗临河，可以眺远。多间茶馆建筑楼台亭舍，养花植木种竹，堆石，所用的杯盘匙筋，无不精美60。茶馆、餐馆还应召到顾客家服务，即主人请二三桌酒，叫「茶酒」之商家前来伺候1。由此亦见餐饮业的发达。

\BookSubheading{（三）美容健身的衣饰与沐浴、理发业}

衣着、首饰、化妆，讲究修饰的人极其在意，即使马虎的人对衣服的质地、颜色也会稍加留心，扬州人是注意美容的典型。

惺庵居士说扬州人「服饰竞时髦」62，李斗则云「扬郡着衣，尚为新样」。扬州府城多子街，因两边多绸缎铺，又名缎子街。在乾隆五十年代以前，流行缎料

\vfill
\begin{BookNotes}
\BookNote{54.}{《扬州画舫录》，第26页。}
\BookNote{55.}{《扬州画舫录·自序》。}
\BookNote{56.}{《扬州风土记略》，第49页。}
\BookNote{57.}{林溥，《扬州西山小志》，第42页，《扬州地方文献丛刊》本，广陵书社，2003年。}
\BookNote{58.}{《邗江三百吟》，第119页。}
\BookNote{59.}{《扬州画舫录》，第198。}
\BookNote{60.}{《扬州画舫录》，第26页、第198页。}
\BookNote{61.}{《邗江三百吟》，第72页。}
\BookNote{62.}{《望江南百调》，第1049页。}
\end{BookNotes}

% 原书第 397 页（PDF 第 410 页）
\BookPage{397}{410}
\noindent 八团花纹，后变为大洋莲；衣服的颜色崇尚三兰、朱、墨、库灰、泥金黄，乾隆五十二年（1787）大将军福康安前往台湾平定林爽文起事，路经扬州，穿的是膏粱红、樱桃红色衣装，人们于是喜爱膏粱红、樱桃红，名曰「福色」63。男子穿的蝴蝶履，为「扬式名鞋」，世传「苏杭人极称羡」，林苏门以为一双鞋有什么金贵的，不相信传闻是真的，及至到杭州始知实情如此，每双鞋价值高达一两二三钱64。女裙，原来将缎料裁成若干条，使用金线对每条镶绣花边，拼结成裙，谓为凤尾；乾隆后期改为百折裙，使用整幅缎料折成细道，常服是二十四折，名玉裙。这些成衣，往往由翠花街成衣店制作65。

翠花街，市肆韶秀，原来是珠翠首饰铺。扬州人常用义髻，有蝴蝶、望月、花篮折项、懒梳头、双飞燕、到枕松、八面观音等，异于他处66。

香料，多用于化妆品，清朝人喜爱佩戴香囊。李斗说「天下香料，莫如扬州」。又是扬州第一。制作的商号，以戴春林为上，张元书次之67。《望江南百调》说到戴家香料店：「扬州，比户戴春林，一样牌题名士手，几番香醉美人心，脂粉旧驰名。」68戴家名声在外，固然香料上好，也是使用者多，扬州人喜用香料制品。既然是「驰名」，必是外地人追求探购而消受。张家也具竞争力，江畹香署理山东巡抚，为乡试监临，令张元书店铺制造汉瓦、奎璧等形的香料，凡乡试诸生各赠送一枚，其后张家名其为「状元香」，士子探购，以图吉利69。

扬州俗谚:「早上皮包水，晚上水包皮。」说的是扬州人的生活习惯:皮包水，指上述的饮早茶；水包皮，则谓泡澡堂。无疑，澡堂业兴旺。《望江南百调》咏曰:「沐浴有跟池，扶掖随身人作杖，摩挲遍体客忘疲，香茗沁心脾。」70是

\vfill
\begin{BookNotes}
\BookNote{63.}{《扬州画舫录》，第194页。}
\BookNote{64.}{《邗江三百吟》，第80页。}
\BookNote{65.}{《扬州画舫录》，第195页。}
\BookNote{66.}{《扬州画舫录》，第195页。}
\BookNote{67.}{《扬州画舫录》，第195页。}
\BookNote{68.}{《望江南百调》，第1045页。}
\BookNote{69.}{《扬州画舫录》，第195页。}
\BookNote{70.}{《望江南百调》，第1048页。}
\end{BookNotes}

% 原书第 398 页（PDF 第 411 页）
\BookPage{398}{411}
\noindent 说洗澡与按摩同时进行，令人身心舒泰。扬州的浴池城厢内外皆有，互相争竞。浴池用白石砌边，区分大中小池，水温不一，供顾客选用。更衣的座位，分设座箱、站箱，是雅间与普通间之别。洗浴之外，有侍者按摩。备有香茶碧酒，令人极尽享乐。沐浴在节庆日、喜庆日与平日不同，除夕沐浴叫「洗邋遢」，端午谓之「百草水」，迎亲前夕沐浴，动辄花费数十两银子71。

理发业者有名为「痒上」的服务，顾客数人聚会，一人作东，令理发师按摩，他不捶而捏，使肩头不痛而痒，主人给钱，客人也给小费72。扬州还有类似今日的足疗，业者对客人的足部按摩，有修、捏、刮、出血诸名目，东台有修脚处，理发师亦兼理73。

\BookSubheading{（四）戏曲重镇}

扬州人戏曲欣赏，眼福不浅。被后人认为，乾嘉时期各地戏曲来扬州串演，令此地成为戏曲演出中心，徽班进京的准备场所，京剧的最初孕育地74。

戏班，两淮盐政例养戏班，以备皇帝巡幸、官员应酬之用，是以扬州戏剧业不会不发达。正如钱咏所说：「梨园演戏，高宗南巡时为最盛，而两淮盐务中尤为绝出。例蓄花雅两部，以备演唱。」他基本上是转述李斗的话。李斗讲，「两淮盐务例蓄花雅两部，以备大戏。」又指出花、雅之别：「雅部即昆山腔；花部为京腔、秦腔、弋阳腔、梆子腔、罗罗腔、二黄调，统谓之乱弹。」76李斗是行家，钱咏即谓其「精于音律」77。他指出在扬州昆曲居主流地位，昆腔班起始是商人徐尚志招苏州名伶组成老徐班，而后有大洪班、德音班，又征花部春台班，演大戏。人们看重昆腔，演出称作堂戏。戏班入城演出，先至梨园总局祷祀，谓之挂牌，次

\vfill
\begin{BookNotes}
\BookNote{71.}{《扬州画舫录》，第26页。}
\BookNote{72.}{《邗江三百吟》，第76页。}
\BookNote{73.}{《扬州风土记略》，第48页。}
\BookNote{74.}{前揭杜瑜弟文，《扬州研究》，第59页。}
\BookNote{75.}{《履园丛话》卷12演戏〉，第332页。}
\BookNote{76.}{《扬州画舫录》，第107页。}
\BookNote{77.}{《履园丛话》卷12〈度曲〉，第331页。}
\end{BookNotes}

% 原书第 399 页（PDF 第 412 页）
\BookPage{399}{412}
\noindent 于司徒庙演唱，名曰挂衣，戏班的优劣，以戏钱多寡为差等，自七两三钱至三两六钱，分为五等，戏班人数之多，常常达百十人。有许多出色演员，小旦马天保，「色艺无双，演占花魁醉归，有娇鸟依人最可怜之致」。秦腔名伶魏长生（三儿）到扬城投奔布政使职衔、总商江春，演戏一出，赠以千金，实为天价。他泛舟湖上，风闻的妓船尽出，争与魏船靠近，为此用船桨相击，而其举止自若，意态苍凉78。

扬州本地戏，属于乱弹类，戏文亦间用元人百种，音节服饰极为土气，谓之草台戏，演出于祷祀场合，名曰台戏，以旦角为正色，配以丑角，搭伙演出79。

专业戏馆，仿自京师南城外戏馆，有固乐园、阳春茶社、丰乐园。观众多，林苏门形容踊跃情形，「呼朋逐队观者如堵，细雨邪风坐稳身」80。戏院演出，不受风雨影响，看客稳坐欣赏剧码。

清唱，以外净老生为大喉咙，生旦词曲为小喉咙，丑末词曲为大小喉咙，刘鲁瞻为小喉咙，为刘派，兼工吹笛。清唱在赛会时举行，争相斗曲，以画船停留多少定胜负81。

曲艺门类颇多。评话，说演义故事，如吴天绪「三国志」，徐广如「东汉」。徐氏在成名之前，苦读汉魏文三年，故能吐属渊博文雅，为士大夫所欣赏。艺人浦琳，艺名皮五，自撰《清风闸故事》，演出中「揣摩一时亡命小家妇女口吻气息」，闻者欢笑，感情随之变化。他还善于说笑话，演口技82。顾汉章的说书，亦负盛名83。大东门书场，中间设书台，观众四面团坐，门悬书招，上书说书人姓名，开讲书目，屋主与说书人分单双日收钱，日收入至一千者为名家。各门街巷皆有书场84。

\vfill
\begin{BookNotes}
\BookNote{78.}{《扬州画舫录》，第122页、第125页、第131页、第132页。}
\BookNote{79.}{《扬州画舫录》，第130页、第133页。}
\BookNote{80.}{《邗江三百吟》，第108页。}
\BookNote{81.}{《扬州画舫录》，第254页。}
\BookNote{82.}{《扬州画舫录》，第205页、第258页。}
\BookNote{83.}{董伟业，《扬州竹枝词》。}
\BookNote{84.}{《扬州画舫录》，第207页。}
\end{BookNotes}

% 原书第 400 页（PDF 第 413 页）
\BookPage{400}{413}
扬州弹词（弦词），与苏州弹词，共为江南弹词的主流85。

打十番鼓，使用十种以上乐器，《望江南百调》云:「扬州好，鼓乐十番多，豪竹分明如法曲，哀思婉转当清歌，缓急应云锣。」86打十番中的「梦香词」、「扬州好，新乐十番佳，消夏园亭雨夹雪，冶春楼阁蝶传花」。讲求对仗87。

小唱，音调变易快，适应观众需求，乾隆中期流行哀戚之声，后来时兴满江红、湘江浪，有的艺人将牡丹亭、占花魁谱为小曲演出88。

口技，井天章善学百鸟声，游人招至画船，与鸟斗鸣。另有艺人名画眉扬，可知技艺之绝89。

傀儡戏。惺庵居士咏道：「傀儡戏登场，凡事由人阴簸弄，此身枉自负昂藏，木偶也冠裳。」90

杂技门类也多。有舞双叉的，张天奇出名1。杂耍，艺人来自四方，有杆戏、饮剑、壁上取火、走索、弄刀、舞盘、风车、菠米、撮戏法、飞水、滚灯等节目92。孙呆周变戏法，技艺出众93。相扑，又名摆架子，唐代相传而来94。各阶层人士皆喜观赏杂技，给人影响深刻，故《望江南百调》留下它的技艺纪录，且艺人来自淮北:「把戏铛姣娆，走马柳腰轻贴地，行绳莲瓣欲凌霄，淮北土音娇。」95

\vfill
\begin{BookNotes}
\BookNote{85.}{参阅严岚：〈悲欢离合胸中记只在三弦一拨间—写在〈扬州弹词·审刁案〉出版之际〉，《古籍新书报》第95期，2010年7月28日。}
\BookNote{86.}{《望江南百调》，第1055页。}
\BookNote{87.}{《扬州画舫录》，第256页。}
\BookNote{88.}{《扬州画舫录》，第257页。}
\BookNote{89.}{《扬州画舫录》，第259页。}
\BookNote{90.}{《望江南百调》，第1055页。}
\BookNote{91.}{《扬州竹枝词》，第9页，《扬州地方文献丛刊》本，广陵书社，2003年。}
\BookNote{92.}{《扬州画舫录》，第264页。}
\BookNote{93.}{《扬州竹枝词》，第8页。}
\BookNote{94.}{《扬州画舫录》，第265页。}
\BookNote{95.}{《望江南百调》，第1055页。}
\end{BookNotes}

% 原书第 401 页（PDF 第 414 页）
\BookPage{401}{414}
\BookSubheading{（五）发达的书院教育}

扬州有很多教育机构，大体上分府州县学、书院和义学三种，书院以设在郡城的安定、梅花和设在仪征的乐仪三间最为著名。扬州府、两淮盐政皆宣导书院建设，盐商出力最多。书院院长皆饱学之士，培养出大量人才，成为扬州是文化城市的重要内涵和标志。

安定书院，系巡盐御史胡文学创办于明朝崇祯十五年（1662），雍正十一年（1733）巡盐御史高斌、运使尹会一，「以广陵名郡，人文渊薮，亟宜振兴，议即旧址重建，谕商公捐」96，商人和士人也希望仿照省会的书院规制办好书院，众商乃捐银修葺一新。所收学生人数先后有所不同，初期定额正课生六十名，后增至百人，各个时期有所变化。学生分正课、附课、随课三类。正课生每年给膏火银36两，附课生12两，住院学习的，于常额外日增膏火三分，此外参加乡试，资送路费，中举者，则予树旗杆、送匾额的荣誉。梅花书院，创办于明中期，由地方官和盐官主持，先后名甘泉山书馆、崇亚书院，明末废。雍正十二年（1734）绅商马曰琯出资重建，所谓「独任其事，减衣节食，鸠材命工…不期月而落成」97，定名梅花书院。乾隆四年（1739）由盐衙支给诸生膏火费，收留学生六十名，亦分正课、附课、随课三种，给膏火费同于安定书院。安定、梅花书院平时有考试，由地方官主持的为官课，由书院掌院进行的为院课，均系每月进行一次，奖励优等者银两，连续获奖者升等。乐仪书院，乾隆三十三年（1768）知县倡建，盐院、盐商支持，正、附、随课生多时达二百四十人，正、附生均享受膏火费，亦有月课奖励银。经费较为充足，乾隆五十九年（1794）至嘉庆五年（1800），每年用银二千二百余两8。《扬州画舫录》说安定、梅花两书院的掌院，「皆知名有道之士」99。安定书院院长全部是进士出身，梅花、乐仪多数院长是进士，间有举人出身的。康雍乾嘉时期安定书院院长有，王步青，翰林院检讨；储大文，翰林院编修；王峻，御史；查祥，编修；陈祖范，国子监司业；邵泰，编修；沈起元，光禄

\vfill
\begin{BookNotes}
\BookNote{96.}{嘉庆《两淮盐法志》卷53〈书院〉，同治九年扬州书局版。}
\BookNote{97.}{光绪《增修甘泉县志》卷16〈学校〉，中国方志丛书本，华中地方第408号第3册第1072-1085页。}
\BookNote{98.}{道光《仪征县志》卷18学校·书院〉，光绪版。}
\BookNote{99.}{《扬州画舫录》，第65页。}
\end{BookNotes}

% 原书第 402 页（PDF 第 415 页）
\BookPage{402}{415}
\noindent 寺卿；刘星炜，侍郎；杭世骏，编修；蒋士铨，编修；吴玨，内阁中书；赵翼，探花、道员；张涛，翰林院侍读；吴锡麒，国子监祭酒。梅花书院院长有姚鼐，刑部郎中；茅元铭，内阁学士；蒋宗海，内阁中书；胡长龄，廷试第一。乐仪书院院长亦为有名人士，有沈廷芳，中乾隆博学鸿词科、按察使；秦黉，编修；王芑孙，乾隆召试举人；赵翼、吴锡麒等人均出任过院长100。这些人中，有的是诗坛领袖，有的是史学大家，有的是桐城派古文集大成者，几乎在《清史列传》、《清史稿》中立有传记，文章被收在《皇朝（清）经世文编》中。书院有明师，又给膏火费，发放岁科两试及乡试路费，颇能吸引学子就学，所谓「来学者，四远麋至」101，「四方来肆业者甚多」102。生源广泛，不限本府本省，更易出人才。三所书院，均以诸生、监生为教育对象，培养出一批科举人才、学者和官员，诚如李斗所说，「能文通艺之士萃于两院（指安定、梅花书院）者极盛」，他们中有管一清，进士，善属文，工诗；杨开鼎，翰林，道员，工书法，有诗集；梁国治，少时肄业于此，官至大学士；谢溶生，刑部侍郎；蒋宗海，前述书院院长；秦黉、秦恩复父子，皆进士、编修；任大椿，及其弟子汪廷珍，均进士，经学家；吴楷，召试中书，工诗文辞赋；段玉裁，有名的古文字学家；李惇，通三礼；王念孙、王引之父子经学大家；汪中，经学大家；刘台拱，精于三礼；韦佩金，尤通于时文，同学称为「文虎」；洪亮吉，博通经史，今世以人口论著称；孙星衍为经学、音韵学家；顾九苞，进士，贯通经史等等103。上面所述基本上是嘉庆及其以前的事，咸丰间扬州各书院毁于战火，同治中恢复安定、梅花以及广陵书院，由盐衙将专门犒赏军队的一部分银子，经过两江总督曾国藩的批准，移作书院的经费。

三所书院之外，扬州还有维扬书院，系明朝中叶巡盐御史雷应龙创建，后废；敬亭书院，两淮商人于康熙二十二年（1683）建立，令士子诵读其中，后废毁。与书院并立的是义学。府城原有义学三间，后增为五所，嘉庆五年（1800)商人洪箴远等，因郡城广大，义学太少，愿意资助，得到盐衙允准，遂于十二门各

\vfill
\begin{BookNotes}
\BookNote{100.}{嘉庆《两淮盐法志》卷53书院〉。}
\BookNote{101.}{嘉庆《两淮盐法志》卷53〈书院〉、卷55〈碑刻〉。}
\BookNote{102.}{《扬州画舫录》，第66页。}
\BookNote{103.}{《扬州画舫录》，第67页。}
\end{BookNotes}

% 原书第 403 页（PDF 第 416 页）
\BookPage{403}{416}
\noindent 设一所104，每学三十人。甘泉县邵伯镇有甘棠义学，康熙间士民公建，嘉庆八年（1803）生员蒋和重修。扬州营义学，乾隆三十八年（1773）建，盐政资助，专课营兵子弟105。

\BookSubheading{（六）文化学术及扬州学派}

笔者仅谈学术方面的扬州学派、刻书藏书、博物鉴赏、琴棋书画与扬州八怪、诗文会五事。

今世学者认为清代学术界产生扬州学派，与惠栋吴学、戴震徽学并称。吴学、徽学被认为是专精之学，而扬州学派是通学106，即为融会贯通之学，吴、徽、扬三学，连同北方的颜李学派、湖南的湘学、经世致用之学，促进了清代学术的发展。扬州学派关注经学的研究。张舜徽指出，扬州学派代表人物之一焦循「于经学造诣甚深，著述多而精」，其子廷琥「亦于穷经为长」，「父子自相师授，有讨论之乐」107。笔者曾就仪征籍学者检索史料，得知在《清史列传》、《清国史》的〈儒林传〉、〈文苑传〉有十一位学者拥有传记，其中许珩的易学研究，被认为是继汉朝人郑玄的《周易》注疏最有成就者108；江昱的痴迷于《尚书》研究，致令袁枚说他是「经痴」109。江藩撰着《汉学师承记》和《宋学渊源记》，论述经学史；阮元主编的《畴人传》，是中国科技史的滥觞之作。均是扬州学派通贯之学的表现。

扬州文化的一个重要特点是这里为刻书重镇与富有藏书家。早在康熙年间，两淮盐政、江甯织造曹寅奉命在扬州编印《全唐诗》，印制《渊鉴类函》、《律吕正义》等巨著，他边刻印边进呈康熙帝御览，满意了，继续雕版，所以是精益求精，刻板技艺高超；他自己另刻有关艺术史方面的书籍十二种11°。马曰琯刻印小学类

\vfill
\begin{BookNotes}
\BookNote{104.}{嘉庆《两淮盐法志》卷53〈书院〉。}
\BookNote{105.}{光绪《增修甘泉县志》卷十六〈学校〉，成文出版社《中国方志丛书》本，华中地方第408号第3册第1077页。}
\BookNote{106.}{参阅前揭王振忠文，《扬州研究》第503页。}
\BookNote{107.}{《清人笔记条辨》，第212页，中华书局1980。}
\BookNote{108.}{《清国史·儒林传下》卷18〈许珩传〉，第12册第640页，中华书局1993年。}
\BookNote{109.}{《清国史·文苑传》卷20〈江昱传〉，第12册第833页。}
\BookNote{110.}{《扬州画舫录》，第51页。}
\end{BookNotes}

% 原书第 404 页（PDF 第 417 页）
\BookPage{404}{417}
\noindent 图籍说文、广韵、玉篇、字鉴等，被称作「马版」111。乾隆四十二年，盐政伊龄阿奉旨于扬州设局修改曲剧，四年事竣，总校黄文赐，分校凌廷堪112。嘉庆帝下令编印《全唐文》，在扬州设淮南书局刊刻，同时刻印《明鉴》113。当地人印书，或盐商为外地学者刻书，是屡见不鲜之事。如马曰琯为朱彝尊梓刻《经义考》114；进士出身的程崟辞官后，「选定明代及本朝古文，次第付梓」；因少年时代与桐城派开创人方苞接近，为其刊刻《望溪全集》115。总商鲍志道，幼时读《论语》、《孟子》，无善本，发迹后细加校正付梓，藏诸家塾。六七十年代扬州有广陵古籍刻印社，线装出版名贵古籍，是其时全国唯一的线装图书出版社，这是继承清代扬州出版事业的传统。

扬州藏书家之富，首推马曰琯、马曰璐兄弟，家有小玲珑山馆藏书楼，乾隆朝编修《四库全书》，由马曰琯子马裕进呈可备探摘的图书七百七十六种，内有多种宋版书，从而获得《古今图书集成》、《平定伊犁御制诗》、《平定金川御制诗》等赐书的殊荣116。藏书多，由目录学著述的出现反映出来。马氏有《丛书楼目录》；秦恩复有《石砚斋书目》；陈木礼有《瓠室》。扬州人富有藏书和献书朝廷，促成乾隆帝在扬州设立文汇阁，收藏四库全书，令扬州人有进一步治学的条件。与刻书、珍藏书画相一致，装裱业发达，装裱古迹珍本、名画技术高超，如匠人叶御夫，「旧画绢地，虽极损至千百片，一入叶手，遂为完物」117。

博物鉴赏，是雅玩之举，也是财富象征，为扬州人所喜爱。林苏门说：「今年讲究人家，未有不学雅而设几方好端砚者。」并举吴绍浣、江藩「识砚」为例118。「未有不学雅」，表明学雅是时尚，是人们的追求。惺庵居士的歌咏表述扬州人的收藏骨董和买卖字画，但买卖者未必识货:「扬州好，古玩富收藏，周鼎汉砖箱什

\vfill
\begin{BookNotes}
\BookNote{111.}{《扬州画舫录》，第88页。}
\BookNote{112.}{《扬州画舫录》，第197页。}
\BookNote{113.}{《扬州访旧录》，第29页。}
\BookNote{114.}{《扬州画舫录》，第88页。}
\BookNote{115.}{嘉庆《江都县续志》卷6〈人物〉，第1册第197页。}
\BookNote{116.}{《扬州画舫录》，第88页。}
\BookNote{117.}{《扬州画舫录》，第207页。}
\BookNote{118.}{《邗江三百吟》，第91页。}
\end{BookNotes}

% 原书第 405 页（PDF 第 418 页）
\BookPage{405}{418}
\noindent 袭，宋书元画锦装潢，癖嗜有萧郎。」「扬州好，古董客高谈，金石刮磨奇货视，丹青装点古人参，真赝几曾谙。」119虽说人们古董知识水准不一，而精于此道者颇有其人。童玨，本身擅长绘画梅花，收藏古今人诗文集，能见到的尽入家门，尤其是「精别古画铜磁玉器金石钱刀」，几乎精通博物学的各门各类120。马曰璐收藏古琴，所藏雷琴，系唐玄宗开元二年（714）制品121。

琴棋书画，为文人应该具备的艺能，其他人要拥有这类本事，不必做文人，但有了文人的名望，故致力于琴棋书画者多。说到绘画，「扬州八怪」为当时及后世人们津津乐道，形成画派，反映画坛的创新精神。郑板桥，工画竹，有八分书与楷书相杂，自成一派，与其唱和者甚多122。希望收藏他的绘画者，不惜金钱:「时人但以字之怪，画之随意，不惜分金而换易之。」123扬州书画家极多，兼之过客往来，代不乏人。如禹之鼎，出入于宋元，自成一家，写真多白描，娟媚古雅124。书画界产生强烈商品意识，尤堪顾及。郑板桥定出「润格」，名曰「笔榜」，对其书画明码标价，以质论价，绝不同求画者讲交情，不怕人家议论文人爱财，理直气壮。其「笔榜」是这样写的：「大幅六两，中幅四两，小幅二两；书条、对联一两；扇子、斗方五钱。」对于讨价还价、以物品顶代银钱的顾客毫不通融，于乾隆二十四年（1759）特做说明：「凡送礼物、食物，总不如白银为妙…送现银则中心喜悦，书画皆佳。礼物既属纠缠，赊欠尤为赖账…画竹多于买竹钱，纸高六尺价三千。任渠话旧论交接，只当秋风过耳边。」125为他人作文，收取润笔，自古有之，买卖字画多有前例，然而公开定出价格，可能是自郑板桥始，在强烈商品意识下，他的勇气令后人赞叹。郑板桥的作风，为清人继承，也是订出润格，名家让顾客先交钱，却常常不及时交画:「先惠笔资悬润格，任追画债懒挥毫，几辈负名

\vfill
\begin{BookNotes}
\BookNote{119.}{《望江南百调》，第1050页。}
\BookNote{120.}{《扬州画舫录》，第195页。}
\BookNote{121.}{《扬州画舫录》，第207页。}
\BookNote{122.}{《扬州画舫录》，第232页。}
\BookNote{123.}{《邗江三百吟》，第98页。}
\BookNote{124.}{《扬州画舫录》，第38页、第40页。}
\BookNote{125.}{《郑板桥全集·板桥》，第243页。}
\end{BookNotes}

% 原书第 406 页（PDF 第 419 页）
\BookPage{406}{419}
\noindent 高。」126在扬州侨寓的画家施胖子，善写真，画美人，长者丈许，小的半寸，取酬三十两银子；杨良，工画驴，画一驴取值牛肉一斤，人称「杨驴子」127。

琴及琴学，徐祎为布政使年希尧刊刻《澄鉴堂琴谱》，徐锦堂着《五音斋琴谱》，并称二徐128。

围棋。国手韩学源，京师授徒，天下无敌129。仪征黄龙士以弈圣著名于清初，乾嘉间僧秋航继起，范西屏以国手寓扬，围棋遂为世重130。他撰有《桃花泉奕谱》，施定庵作《奕理指归》等，先后辉映131。

扬州特色的诗文会，由具有高文化素养绅商主持并出资，招来文人雅士和幕宾，研究学术和文艺，举办歌咏宴会，出版诗集。这正是扬州作为文化城市的一种标志。「扬州诗文之会，以马氏小玲珑山馆、程氏筿园及郑氏休园为最盛」。小玲珑山馆是马曰琯、马曰璐兄弟的，前述藏书事已有所道及。马曰琯为诸生，「好学博古，考校文艺，评骘史传，旁逮金石文字」，着有《沙河逸老诗集》，本身就是一位力学的学者，康熙帝南巡，两次赐予御书、满人食品的克食，进宫为皇太后祝寿。马曰璐，工诗，著作《南斋集》，不乐仕宦，因而拒绝博学宏词科的召试，兄弟二人被人称为「扬州二马」。「小玲珑山馆」，拥有两栋藏书楼。他们热情接待士人，研讨诗文，「所与游皆当世名家，四方之士过之，适馆授餐，终身无倦色」132。如杭州人厉鹗诗文颇受名家欣赏，「搜奇嗜博，馆于扬州马曰琯小玲珑山馆者数年，肆意探讨，所见宋人集最多，而又求之诗话、说部、山经、地制，为《宋诗纪事》一百卷、《南宋院画录》八卷」。马曰琯对他多方照顾，因六十岁尚未生子，特辟住宅为他纳妾。及至他回乡亡故，消息传来，又为他设灵位祭奠133。

\vfill
\begin{BookNotes}
\BookNote{126.}{《望江南百调》，第1050页。}
\BookNote{127.}{《扬州画舫录》，第49页。}
\BookNote{128.}{《扬州画舫录》，第207页。}
\BookNote{129.}{《邗江三百吟》，第95页。}
\BookNote{130.}{徐谦芳，《扬州风土记略》，第49页。}
\BookNote{131.}{《扬州画舫录》，第260页，钱咏在《履园丛话》不以他们为然，见第314页。}
\BookNote{132.}{《扬州画舫录》，第86页。}
\BookNote{133.}{《清史列传》卷71〈厉鹗传〉，第18册第5833页，中华书局1987年。}
\end{BookNotes}

% 原书第 407 页（PDF 第 420 页）
\BookPage{407}{420}
\noindent 郑氏出忠义人士，三度撰写《休园志》。总商江春，诸生出身，工制艺，精于诗，迎接乾隆帝南巡，报效甚多，得到君主欢心，出席千叟宴，赐布政使秩衔，是所谓布衣而与天子交游者。他广结宾客，建随月读书楼，请人选时文付梓行世，名《随月读书楼时文》。他死后，每日来祭灵而不报姓名的有十多人，所以有人将他比作孟尝君。诗会，先订会期，届时，在园中设置若干条案，每案供一人使用，案上放置纸笔墨砚和水注，诗韵，茶壶茶碗，菓合、茶食合。诗牌，是象牙制品，一寸半见方，每人分得数十字，或百余字，凑集成诗，最难工妙。与会者诗成，即行刻印，次日送各人审阅，在三日内若有改动，则重新刻印。马曰琯为此刊有《韩江雅集》。每会，备有极为珍美食品；听曲，先在陋室听老艺人奏曲，次后在豪华露厅观赏妙龄男女艺人演奏。所作之诗，至乾隆末年仍在流传，有汪中的:「叶脱辞穷巷，莲衰埽半湖。」张四科的:「舟棹恐随风引去，楼台疑是气嘘成。」等佳句134。

\BookSubheading{（七）宗教文化与信仰}

扬州人的宗教和神灵信仰是多元的，各种各样，与其他地方多有相同，但有其特点，即与皇帝关系密切，从而富有特定的文化内涵；宗教的多元性，反映扬州作为水道码头，尤其是历史上的内河国际海港的文化包容性；强调伦理传统性与时代性的结合。

扬州的佛道及民间信仰诸神寺宇，多有同清朝康雍干三帝密切关系者。这一层因缘，为它们增添浓重的色彩，更其出名，和富有文化内涵。比如栖灵寺，隋代创建的古刹，康熙帝于四十四年（1705）赐「澄旷」匾，乾隆三十年（1765）赐名法净寺。是以顾銮在《广陵览古》中说该寺:「清代增建僧寮、殿阁，恭迎列圣临幸。寺本名刹，得此亦壮烟雨楼台之景色。」135为康熙帝、乾隆帝南巡接驾，盐衙、地方官和盐商对原有寺宇大肆修缮扩建，为皇帝驻跸和游幸。高旻寺，康熙三十八年（1699）赐匾额，四十二年（1703）准许盐商建行殿，康熙帝临幸，赐御

\vfill
\begin{BookNotes}
\BookNote{134.}{《扬州画舫录》，第180页。}
\BookNote{135.}{顾銮，《广陵览古》，第34页。}
\end{BookNotes}

% 原书第 408 页（PDF 第 421 页）
\BookPage{408}{421}
\noindent 书匾，撤大内所供金佛于寺，四十四年（1705）、四十六年（1707）赐诗文法物；雍正间重修、购置田业；乾隆帝驻跸于此136。天宁寺，康熙帝赐匾额四、楹联二。乾隆帝赐匾额七、楹联八，一方匾曰「省方设教」，云其南巡在于了解民情，施行教化。而康熙帝的一方匾曰「般若妙源」，寄望寺院弘扬佛法137。法海寺，康熙帝二十三年（1684）两度临幸138。三汊河高旻寺，康熙四十三年（1703）淮商加建行殿，供皇帝巡幸之用。观音山寺，乾隆二十一年（1756）商人建造，以备乾隆帝南巡休憩139。扬州道家罗荣光，自号「寻真子」，康熙三十一年（1692）、三十八年（1699）两度应召进京，康熙帝赐书「通幽索隐」，他在扬州建斗姆宫为康熙帝祝釐，四十六年（1707）南巡，临幸其宫，御书「大智光」额。雍正三年（1725）雍正帝赐给他银二千两140。三汊河关帝庙，乾隆历次南巡，众商修葺141。传云欧阳修所建平山堂，成为后世士绅活动场所，屡加修葺。康熙帝两度临幸，撰有〈御制平山堂诗〉，少詹事高士奇和诗；康熙帝因扬州知府尹会一撰〈重修平山堂记〉，赐平山堂「贤守清风」额，宣导地方官廉洁；又书「怡情」赐禅师142。邵伯镇露筋祠，乾隆四十五年（1780）、四十九年（1784）乾隆帝两次南巡过此，两淮商人重葺。宝应神龙祠:乾隆二十六年（1761）建造，乾隆历次南巡淮商捐资修葺。宝应湖神庙，乾隆三十年（1765）修建，四十五年（1780）、四十九年（1784）乾隆帝南巡，淮商先后修缮143。

唐宋元时期，扬州有数以千计的阿拉伯人、波斯人、伊斯兰教回回人，并有回回人担任地方长官的。宋元之际的宝合丁（普哈丁）建立仙鹤寺，遗址在扬州南门大街，他的墓地在东关城外老运河东高岗上，俗称回回堂，墓园有元代回回人的墓碑四通，碑文反映回回人的官员、教长、商人、学士、妇女、平民的活动，基本上

\vfill
\begin{BookNotes}
\BookNote{136.}{嘉庆《两淮盐法志》卷52〈杂纪一·祠庙〉。}
\BookNote{137.}{《扬州画舫录》，第84页。}
\BookNote{138.}{《平山堂揽胜志》，第53页。}
\BookNote{139.}{嘉庆《两淮盐法志》卷52〈杂纪一·祠庙〉。}
\BookNote{140.}{嘉庆《江都县志》卷6人物〉，第1册第240页。}
\BookNote{141.}{嘉庆《两淮盐法志》卷52杂纪一·祠庙〉。}
\BookNote{142.}{《平山堂揽胜志》，第64页、第72页、第137页。}
\BookNote{143.}{嘉庆《两淮盐法志》卷52〈杂纪一·祠庙〉。}
\end{BookNotes}

% 原书第 409 页（PDF 第 422 页）
\BookPage{409}{422}
\noindent 反映了元代来华回回人各方面人士的历史片段144。十四世纪上半叶，在扬州有基督教派的圣方济各派教堂一所，景教教堂三所。元代扬州，有基督教堂、伊斯兰教堂与墓园，并有遗迹、遗址保存至今，表明当地有基督教徒、伊斯兰教徒的活动，是多元文化地区，有回回人的后裔定居扬州，扬州名菜张四回子全羊即为明证，说明后世的汉人、满人接纳他们的文化。作为水陆码头，扬州文化具有包容性，也才成为文化名城。

扬州作为商城、水路码头，与同类城镇有一些共同的信仰祠宇，如有财神庙、水神庙、天后宫、金龙四大王庙、江神庙、河神庙。祈求金钱和平安的同时，讲求主流意识的忠孝节义伦理，崇尚名贤，对董仲舒、欧阳修、范仲淹、文天祥、李庭芝、史可法表示特殊的敬意。有董子祠，已如前述。欧阳文忠祠，纪念欧阳修，雍正十年（1730）商人修缮，乾隆十五年（1751）众商重修。范文正祠，春秋二祭。双忠祠，祭祀南宋守扬州的李庭芝和姜才，雍正十二年（1734）年马曰琯修缮。大忠祠，祀文天祥。五贤祠，纪念宋人、明人，雍正十二年（1746）盐政允许众商捐修。萧孝子祠，祀康熙六年（1667）为病母而死的本地人，雍正十二年（1734)马曰琯重建。五烈祠，雍正十一年（1733）甘泉知县改建。对于明清之际史可法的殉节扬州，当地人和外地人均表示景仰，惺庵居士《望江南百调》云:「阁部史公坟，孝子萧祠相接壤，忠臣信国是前身，梅岭仰清芬。」145浙江人姜宸英的〈游平山堂感事有作二首〉，其二咏怀史可法:「朝游城北暮城东，相国名犹满域中。…荒冢至今闻夜哭，可怜弦管醉春风。」146康熙间，在政治高压下，不可能出现众人纷往史可法墓凭吊的场面，但汉人内心深处景仰之。在清朝实行歧视汉人的政策情形下，扬州人纪念文天祥、李庭芝、史可法，表现出汉人潜在的汉民族气节观念的强烈。

\vfill
\begin{BookNotes}
\BookNote{144.}{此处节写杨志玖先生〈元代扬州的回回人〉，文载《扬州研究》。}
\BookNote{145.}{《望江南百调》，第1041页。}
\BookNote{146.}{《平山堂揽胜志》，第125页。}
\end{BookNotes}

% 原书第 410 页（PDF 第 423 页）
\BookPage{410}{423}
\BookSubheading{（八）市政建设与社会救济事业}

清代扬州府城的建设颇可称述，主要是河渠道路的修缮，关注交通的便利与安全；注意防火；收容贫病民人；施钱施药；赈灾。

\BookSubheading{(1)水陆道路的维修与安全保护}

《望江南百调》说得好，扬州是「街小傍河干」147，滨河的城市河流水道的疏通是为要务，扬州又是以水路运输为生命线，地方官、盐官和盐商更加关注。疏浚城河，不断进行。府城水沟易于淤塞，乾隆二年（1737）淮南总商创意修浚，绅商马曰琯自任其住宅所在的广储门至便益门段的疏通，其余十四段众商公修。乾隆二十四年（1759），护城河和城内市河严重淤塞，众商情愿出银一万七千六百两挑河疏通，竣工后，为不令杂物堵塞河道，设专船清除垃圾。城内官井亦相继疏浚。乾隆五十七年（1792）总商洪箴远等公捐重浚。修筑桥梁街道，郡城西北孔道的扬州古雷塘，嘉庆三年（1798）建造石桥，以便行旅；乾隆中，淮商罗琦重修东关大街和城外石码头；总商鲍志道重修新城街道。

为便利水道交通，设立义渡。有的地方不便造桥，为了行人的方便，设立义渡，备置船只和船工，无偿载渡行人。义渡每年的费用，由商人捐助，在江都和仪征有十余个，如徐宁门外二严庵万松义渡，即由商人汪勤裕捐建；钞关门外义渡始由诸生耿兆组捐田供费用，后来耿氏子孙卖田，致使费用无着，嘉庆八年（1803)改归盐务支银。

为维护水运安全，扬州商人设置救生红船，救助遇难船只和人员。船舶设于江都史家港、双港口、大沙洲、大江镇、瓜洲江口、瓜洲江神庙、仪征天池、仪征沙漫洲、金山、焦山、镇江避风馆、高邮甓社湖等处。救生船是大帆船，要经常维修，水手要有工钱，均由众商渐次修补，并不断增加船只，水手工食银亦有定额。

\BookSubheading{(2)防火}

扬州府城人烟稠密，会有回禄之灾，为了救火，城厢设置水仓，备有水炮、水

\vfill
\begin{BookNotes}
\BookNote{147.}{《望江南百调》，第1041页。}
\end{BookNotes}

% 原书第 411 页（PDF 第 424 页）
\BookPage{411}{424}
\noindent 缸及救生用具，并动用水兵服役，所有「修理器具及各处救火兵役，岁需银两，皆出商捐」148。也就是说商人承担了城市消防救火的费用。

\BookSubheading{(3)收容贫民、弃婴的善堂}

有普济堂、育婴堂、收容所等类型。

扬州普济堂，收养无告老民、贫民，康熙三十九年（1700）创立，乾隆九年(1744）因经费不敷，遂于商捐项内每年拨银。瓜洲普济堂，雍正二年（1724)江都贡生耿兆组捐建，随后知府陈宏谋捐俸及募金扩建，耿氏陆续捐田三百亩、银万两为维持经费，雍正九年（1731）盐运使改令淮商经管，每年给银一千二百两，乾隆中岁支银二千五百八十两。邵伯镇有同善堂，即普济堂，知府尹会一建，雍正十一年（1733）盐政高斌岁给银四百两作为经费。

扬州育婴堂，明代末年毁于兵燹，清初商人创其事，顺治十二年（1655）建成，每年需银二千两，郡绅李宗孔暨商人闵世璋、郑元化等捐助，后来难于为继，康熙五十年（1711）运使李陈常接受商人请求，月给银百两。盐政佥派商人管理，时间一久，经管商人以有赔垫，请求退出，运使卢见曾乃于乾隆二十年（1755）令众商议购菜田，作为育婴堂的固定资产和收入，旋于淮南二十八总商内派出二人经理其事，乾隆五十六年（1791）盐政全德将两淮归公田房统交商人承管，补充育婴堂经费。乾隆六十年（1795）盐政扩建乳婴室，盐官与「商人之有力者皆乐输」，由知县和总商董其事，建有乳妇室二百四十间，设妇头，有幼儿医生、外科医生各一人。瓜洲育婴堂，因经费不足，众商于乾隆十七年（1752）公捐银一千二百两。

扬州收养所，乾隆三十二年（1757）建，每年冬季收养冻馁无依靠者，需银五六千两149。

\BookSubheading{(4)施钱施药}

日常的施钱施药，重灾时赈济。

\vfill
\begin{BookNotes}
\BookNote{148.}{嘉庆《两淮盐法志》卷55〈杂纪·普济堂、育婴堂〉；林苏门《邗江三百吟》，第26页。}
\BookNote{149.}{嘉庆《两淮盐法志》卷56〈杂纪·普济堂、育婴堂〉。}
\end{BookNotes}

% 原书第 412 页（PDF 第 425 页）
\BookPage{412}{425}
扬州施药局，盐衙主管，雍正七年（1729）总商黄光德等公捐设立，于每年的五月至八月施舍丸散。殷实之家在城内外设公所一药局，延医，布施丸散。乾隆二十一年（1756）大疫，盐政令商人江助周等增设一个药局，请医生救治，数月后瘟疫消灭，裁撤150。

恤嫠局，嘉庆十一年成立，资助嫠妇每月四百文。由官员捐俸，绅董经理151。

施茶，在各城门，热天为肩挑负贩者设置免费茶水152。

赈灾，乾隆三十六年（1671）维扬被灾，众商捐银二万二千六百七十两买米赈济，在扬州城外设立四个粥厂，每日煮粥，约赈男女四万五千人；泰州、兴化、高邮等处则每日发米数百石，又给灾民棉衣一万件。捐银在五百两以上的，清朝政府给予九品顶带，陈恒升等八人获得这种荣誉。康熙十八年（1679)旱灾，众商捐银三万三千余两，救济饥民八百一十六万余。乾隆三年（1738）扬州旱灾，众商捐银十二万七千两，独捐四万七千两救灾。六年（1741）秋雨过多，淮南商人黄仁德等设厂煮粥赈济两月；次年维扬水灾，捐银六万，黄仁德等公捐二十四万两；乾隆十一年（1746）淮南众商程可正等因灾公捐二十万两；乾隆十八年（1753）两淮灾，商人捐助三十万两，等等153。

\BookSubheading{(5) 盐义仓}

储存粮食，为赈济灶丁之用。扬州东关附近一仓，雍正四年（1726）建立，雍正十二年（1734）商人捐银二十万两购买仓米，至乾隆十一年（1746）贮存粮食六万石，佥商二人管理；广储门外三仓，亦是雍正四年（1726）建，乾隆十一年（1746）储粮十八万石，商人管理。所佥商人，是盐政认可的「老成殷实商人154。

\vfill
\begin{BookNotes}
\BookNote{150.}{嘉庆《两淮盐法志》卷56〈杂纪·药局〉；《邗江三百吟》，第24页；《扬州画舫录》，第207页。}
\BookNote{151.}{《邗江三百吟》，第22页。}
\BookNote{152.}{《邗江三百吟》，第24页。}
\BookNote{153.}{嘉庆《两淮盐法志》卷42捐输三·灾济〉。}
\BookNote{154.}{嘉庆《两淮盐法志》卷41〈恤灶〉。}
\end{BookNotes}

% 原书第 413 页（PDF 第 426 页）
\BookPage{413}{426}
\BookSubheading{(6) 义冢}

雍正十二年（1734），商人黄仁德等奉两江总督之命，捐资于四郊买地十六处作义冢；商人置义冢于徐宁门外；乾隆二十四年（1759）商人黄源德等捐资，分别于南门外、北门外购置义地；乾隆四十三年（1778）商人江春等捐资买西山空地七十亩作义冢，二十亩给僧人耕种纳粮，以维持义冢155。

陈向山捐地设义冢，庵僧经理。

拾骨坛，阮元出资建造，令秋雨庵僧经理156。

市政建设和慈善事业，扬州的特点是具有普遍性和相对的稳定性，居民受益。城市道路、水路交通、消防、善堂等等的基本建设，临灾的救助，既有组织机构的保障，又有活钱的调用，所以提供了居民生活的某种方便，同时因救灾能力较强，以利于民众的摆脱灾难。不能不说，在那时，扬州人的生活品质，可能高于苏州以外地区的人。普济堂之类的建设，大抵是皇帝号召，尽职的、当地有条件的地方官奉行，多因没有或缺少固定田产，不能维持而败坏；扬州的不然，基本上能够坚持操办，乃因绅商不断赞助，盐衙监管有力，如救生红船由商人开始做起，是淮商吴自牧设救生船于金山，又雇渔舟协助救生；方如珽设救生船于避风馆，悬赏格，每救一活人，奖银一两，捞获死亡者给一半，并营建墓地一区，安葬遇难者；汪文学设救生船于燕子矶，另置田百余亩，充常年经费；在镇江、焦山、瓜洲，也立赏格。嘉庆《两淮盐法志》说救生红船就是仿照这些事例建立的：「救生船故事，后所设立，盖仿于此。」157对于盐商的这种作用，巡漕御史陈科捷在讲到扬州育婴堂的建设时说到两种因素：「当道诸公爱人励俗之意，商人乐义奉公之举。」肯定了商人的义举。毫无疑问，扬州盐商协助盐务衙门和地方政府开展社会公共事务、公益和福利事业的活动，颇着成效，提高了社会救灾能力。地方公益和社会救济事业的从事者，在盐商、盐衙之外，土着也在做，前述瓜洲普济堂的肇兴者耿兆组就是江都土着，他还帮助贫民交纳拖欠的钱粮，捐田资助监狱犯人的伙食，设立义塾，

\vfill
\begin{BookNotes}
\BookNote{155.}{嘉庆《两淮盐法志》卷56〈杂纪·义冢〉。}
\BookNote{156.}{《邗江三百吟》，第23页。}
\BookNote{157.}{嘉庆《两淮盐法志》卷56〈杂纪七〉。}
\end{BookNotes}

% 原书第 414 页（PDF 第 427 页）
\BookPage{414}{427}
\noindent 教育乡邻子弟。其兄兆绅，「倜傥好施与」，设义渡，置义冢，荒年出粟赈饥158。晏德宝，世居仪征南门外，濒临大江，附近是淮南制盐所，船舶往来如织，每遇风浪，多有覆溺，乃捐资设立救生船局，乡里称为「善人」；其子行忠，夙承父训，乐于助人，料理仪征所设立的善堂同仁堂，仿照苏州同仁堂规则，严定章程，所司之事，历久不懈159。绅、商合作进行地区的建设，扬州的城市建设和社会公益事业才得以发展。

\BookSubheading{（九）狭邪业的娼妓}

扬州娼妓、变童、赌博之风甚盛。《望江南百调》咏:「扬州好，花窟足勾留，罗袜凌波夸石首，玉箫吹月试珠喉，狎客乐忘忧。」160李斗记叙是，妓馆「每夕燃灯数万，粉黛绮罗甲天下」，有专门教曲的艺人，如邬抡元、方张仙161。道出妓业兴盛的概况。钱咏说:「士大夫俱尚豪华，而尤喜狭邪之游。在江宁则秦淮河上，在苏州则虎丘山塘，在扬州则天宁门外之平山堂，画船箫鼓，殆无虚日。」162不幸，扬州在这方面也出名。

附带说一下，扬州客舍业发达，其客源是「清客、评话、戏法、女班及妓馆母家来访者」163。亦反映扬州城市的特点。

\BookSubheading{三、历史启示：建设文化教育休闲城市}

上面从九个方面了解盛清时期扬州人生活的行为与习尚，就中有积极健康的生活因素，也有消极的成分，今天对它的认知，不难发现它对后人的历史价值，那些有益的成分会给扬州市政建设、扬州人健康生活带来福祉。

\vfill
\begin{BookNotes}
\BookNote{158.}{乾隆《江都县志》卷22。}
\BookNote{159.}{同治《续纂扬州府志》卷14。}
\BookNote{160.}{《望江南百调》，第1056页。}
\BookNote{161.}{《扬州画舫录》，第197页、第199页。}
\BookNote{162.}{《履园丛话》卷7〈醉乡〉，第193页。}
\BookNote{163.}{《扬州画舫录》，第203页。}
\end{BookNotes}

% 原书第 415 页（PDF 第 428 页）
\BookPage{415}{428}
\BookSubheading{（一）奢华之城市不能持久，扬州的盛清光辉难再}

扬州「甲天下」名目多，标志着扬州人高水准的消费和高水准的生活，这高水准有正常的、正当的一面，亦有不正当的方面，则是奢靡成风。吴锡麒说他「往来维扬，见其俗奢靡欲折之以道，而未能也」。他认为奢侈的状况是，「农不勤亩，妇不织机，百金之花，千钱之鲋，家无盖藏，而费乃不赀」164。即不知勤俭而过奢华生活的人具有普遍性，所以势态严重，亟须改变。奢靡者主要是盐业富商，李斗在《扬州画舫录》鲍志道条写尽盐商之奢靡无度，所谓「扬州盐务，尽尚奢丽，一婚嫁丧葬，堂舍饮食，衣服舆马，动辄费数十万」165。虽然鲍志道等人努力改变这种风习有一时之收效，而迎接乾隆南巡又奢华不已。「长夜欢娱日出眠，扬州自古无清画。」166不少富商就是这样地生活。富贾如此，一般人家也不示弱，诸如无论贫富爱戴鲜花，小户人家办喜事也叫评词、弦词艺人来伺候一日；稍微有一点条件的人「学雅」，收藏古董。这就是吴锡麒痛心疾首之所在。反对奢华的士大夫对此无能为力，官府有一定的力量能够有所控制，但不能经久维持，如禁止倡业令下，娼妓「辄生死逃亡不知所之」167。禁令一松懈，依然如故。

淮盐税利在中央政府财政中占据重要地位。嘉庆间监修《扬州府志》的两淮盐政阿克当阿，谓扬州「盐荚之利，邦赋攸赖」168。表明盐利不仅对扬州，更对国家财政有重大意义。他的同僚德庆则云:「东南三大政，曰漕，曰盐，曰河，广陵本盐荚要区，北距河淮，乃转输之咽吭，实兼三者之难，其视江南北他郡尤雄剧。」169说明了扬州在盐政、漕政、河政中的重要地位。扬州的繁荣，离不开盐利和盐政，这是人们的共识。盐利之巨，成为盐商奢靡生活的基本条件，暴富者仿效上流社会生活方式一学雅，各方面消费激增；食盐系官方控制销售，皇帝南巡享受盐商的报效，视为当然，促成盐商的奢华。华靡成风，下层社会仿效，于是愈演愈

\vfill
\begin{BookNotes}
\BookNote{164.}{同治《续纂扬州府志》卷23〈艺文·广陵赋〉。}
\BookNote{165.}{《扬州画舫录》，第148页。}
\BookNote{166.}{《郑板桥全集》，第76页。}
\BookNote{167.}{《扬州画舫录》，第198页。}
\BookNote{168.}{嘉庆《重修扬州府志》，阿克当阿〈序〉。}
\BookNote{169.}{嘉庆《重修扬州府志》，德庆〈序〉。}
\end{BookNotes}

% 原书第 416 页（PDF 第 429 页）
\BookPage{416}{429}
\noindent 烈。

经济条件变坏，奢华生活自然难以为继。历史上扬州之兴盛在盐业和河运、海运，衰败也在海运、运河和盐业。所谓海运，是唐代内河国际港口，有对外贸易的海运之利，而漕粮由运河改道东海、黄海、渤海航运，后世扬州不仅没有海运之利，且受其不利影响；运河输送漕粮，为设在北京的中央政府之命脉，近代以来铁路建设，运河地位一落千丈，现今若运河全线通航，虽可运输一些物资，与昔日之通漕要津已不是一回事；盐税是清朝财政收入大宗，仅次于农业税、商税，位居第三，而两淮盐运司的淮盐之利又居各盐运司之首，后世盐业在国民经济、财政地位下降，供应基地的扬州，先前的地位已不可得。总之，扬州的内河国际港口地位已不可复，运河可复，然铁路、航空、海运，使得运河失落大动脉地位；漕运、盐利在政府与民间经济中地位与明清不可同日而语，不可能大复兴，因而扬州不再可能有盛清时代的全国性经济文化重镇的光辉地位。那么扬州可以再创辉煌吗？应当可以，请看：

\begin{BookQuote}
（二）借鉴生活行为中的高雅文明，建设文化休闲城市
\end{BookQuote}

盛清时代扬州有那么多「甲天下」的行业、活动场所及其文化观念，能不能继承，又如何继承？

扬州兴盛繁荣由徽商、西商（山陕商人）和土着共同创造，徽商，被今日学者认为有「贾而好儒」的特征，他们看重文化，尊重文化，努力成为儒商，像马曰琯、程梦星、江春等人，以及他们的子孙族人，多成为文士，中举、中进士、成学者。商人在暴富之初，固然是学雅一附庸风雅，然而很快提高文化，令生活行为赋有文化品味，并成为自身内在的追求。诗会的举行，是文化生活情趣的突出表现；生活在园林之中，感受人生的自然之乐；观赏文艺表演，陶冶情操；古董鉴赏，是艺能，也是高境界的艺术性生活；鲜花的爱好，是美的享受；工艺品的运用于日常生活，令生活处处洋溢着美感。这些都是高素养的人生活中雅趣的表现，是高尚文化的体现。这种富有高雅文化内涵的生活方式，为古往今来人们所向往，盛清时代的扬州人能够实现，实在是幸运之至。在我国社会处于转型期的当代，粗俗文化流


% 原书第 417 页（PDF 第 430 页）
\BookPage{417}{430}
\noindent 行，虽然有其必然性和某种合理性，但是粗俗文化总是需要提高的，物质生活越富裕，文化水准也应该相应提高，因此盛清时代扬州人的追求高雅生活的精神，是值得今人借鉴的精神文化遗产。

由此笔者想到，扬州盛清时代的经济重镇地位虽不可复活，文化教育名城（不是历史文化名城）是可以建设的。需知没有雄厚经济基础的文化名城是可以出现的，可以独立存在的，民国时代的北平（北京）就是一个典型。

具体地说，建设文化休闲城市，其内容，我想到的是：

图书出版事业。图书出版、销售中心，尤其是线装古籍图书的梓行；建设图书出版博物馆。

园林业。恢复某些私家园林，对它作出高文化品味的介绍，使之真正成为大众休憩场所，而不是高收费、乱收费之地，或高收费的「儿童乐园」；民居的园林化，民众的自身家庭建设。

花卉业。生产、供应季节性鲜花，每日不同时辰的鲜花；培植琼花；制造盆景。我在桂林，早上见到沿街出售鲜花者，年轻女性买了即佩戴，扬州不知是否有此情形？要让鲜花供销全国，乃至世界，如同荷兰阿姆斯特丹之鲜花业。

美容健身业。发展理发业、沐浴业、按摩业、香料业；现代家庭洗浴方便，似乎沐浴业不能存在，不尽然，日本人崇尚温泉浴，浴业同样存在；建设美容业博物馆；这是健康的，而非色情的行业。

手工艺的玉雕、漆器，是艺术品，也是实用品、日用品，如香料容器；精致的与大众化的同时生产销售，不必偏废。

餐饮业。大餐、廉价的小吃并行，需知价廉菜肴，给人亲切感，2009年我的一位同事全家五人到扬州旅行，一餐吃十来样饭菜，才花五十元，很高兴，抒情诗云「共和（饭馆名）之唇宴，小吃数不清。色香味俱好，五人五千分」。高雅的茶馆，餐前餐后均宜，雅致环境享受；餐饮体现出食文化的高雅水准。


% 原书第 418 页（PDF 第 431 页）
\BookPage{418}{431}
博物馆。古器物陈列馆，专题的，综合的，遗址发掘的。

寺庙。佛家文化的见证。

旅游业。各种档次的旅社，以普通为主，不追求五星六星级，高档次；让游客盘亘数日，感到舒适，花费与所值相当，打造最宜人的旅游城市形象；与皇帝有关的实物及遗存、逸闻趣事之展览、介绍，需要警惕帝王崇拜文化心态的泛滥。

我基本上生活在书斋中，务实知识少，以上仅从扬州史资料和几次来扬州开会的印象写出的粗糙建议。我理想的文化休闲城市，对本地居民讲是生活在最宜居的城市，市民文明礼貌，气质高雅，令人起敬；对游客，需要让人置身于优雅文明的环境，实现这样两个目标：

一是散心，充分休息，达到身心愉悦。是积极休闲，令人深切感受到生活的乐趣，对生活有寄托，有盼头，下次再来扬州，或到其他地方旅行。总之，增强人的生命活力。

二是以扬州高素质的文化生活氛围，感染、提高游人的行为修养，做有教养的人，用老话说是做绅士型的人。回顾南北朝时代的士族生活，他们讲求修养，诗书传家，书法、棋品、弹琴、歌唱、舞蹈，样样来得，以至于精通，衣着打扮，无不讲究，甚至男子穿着香料熏过的衣服，刮面，搽红粉，乘坐有长辕的车，脚穿高靴，坐在绫罗座垫上，倚在多种颜色丝线织成的靠枕上，从容出入，人们看他简直像神仙那样神采飘逸170。这种优雅闲适风度是士族的必备条件。士人讲究穿着、仪容、气度、语言、饮食，否则就不配做士族。这就叫「高雅」！要让人有高雅气质，必须加以培养。扬州打造文化休闲城市，给人创造优雅生活环境、气氛，去感染人。这是多么崇高的事业，何不为之！近日在网路上见到章诒和回答记者问，她说：「到柳亚子家去看，那就是诗的天地；到张伯驹家，那就是艺术的天地。」「我见过真正的文人和艺术家是什么样的。人家玩都能玩出贡献来，张伯驹所有的贡献都是玩出来的，他无职业，是散仙，一个人从早上起来就晃悠，晃悠出名堂。老说人家钱多，现在的富翁能玩出张伯驹那样吗？」生活上，「罗隆基先生是比较

\vfill
\begin{BookNotes}
\BookNote{170.}{《颜氏家训·勉学》。}
\end{BookNotes}

% 原书第 419 页（PDF 第 432 页）
\BookPage{419}{432}
\noindent 讲究的，康同璧先生也是比较讲究的，马连良啊程砚秋啊梅兰芳啊言慧珠啊叶氏兄弟啊那都是非常讲究的」171。她主张讲究的高雅的文化，鄙视粗俗文化。现在人们倡言强国、大国，靠什么成为大国、强国？经济、军事实力固然重要，不过不可忽视的是文化实力，即所谓软实力，文化观念、生活方式能不能输出，能不能为他人接受才是重要的。粗俗的国民，没有教养的人，能让人尊敬吗？说白话，能让人服气吗？所以建设文化休闲城市，提高民众文化气质，意义可不一般。另外，作为个人讲，「人生无常」，「人生几何」，因此及时行乐就有一定道理，这是不应讳言的。当然，乐有雅乐和淫乐之别，笔者所讲的自然是前者的雅乐。

讲文化休闲城市，是不是享乐主义、高消费主义？明代泰州学派讲「百姓日用即道」，关注民间生活。这种生活观念难道不能传承！享乐，正当的、正常的享乐，是理所当然的，是人的生理、心理所需要的，特别是在现代社会，人们生活节奏快，心情紧张，更需要休闲调节，理应享受，故而旅游业才在世界范围内得到快速发展。享受需要有序、有度，是正当之乐，而不能淫乐，扬州一些富商生活的「无清书」，颠倒白天黑夜，是生活无度无序，与正当享乐不是一回事，自然要引以为戒。高雅生活，是否高消费？两者有关联，但绝非一事。高雅生活自然要有花费，然而丰俭在于个人掌握。前述餐饮、旅社之档次宜多，就是为给消费者提供多种选择的机会，令各种需求的顾客满意。再说高消费不是乱消费，乱花钱是不会花钱，是暴发户的粗鄙行为，自卑心理的反映，令人齿冷、鄙视。正是他们亟须提高社会文化素质。

\vfill
\begin{BookNotes}
\BookNote{171.}{章诒和，绝不会出卖我的灵魂〉，见于网路。}
\end{BookNotes}

% 原书第 420 页（PDF 第 433 页）
\BookPage{420}{433}
\BookSubheading{（三）发扬尊重教育事业精神，宣导家教和办好高等学校}

扬州本是文化之乡，民间历来看重文化教育。扬州人崇祀董仲舒表明了这种心愿和理想。《望江南百调》就董子祠咏道:「董子有崇祠，『正谊』额犹留睿藻，大儒坊合表遗徽，谁更下书帷？!」172「谁更下书帷」，董仲舒「三年不窥园」173，专心学业，他的榜样，鞭策青少年学习文化，哪敢离开书屋贪玩，如此好学，自然成才。扬州人的家庭教育，以学业为重，郑板桥咏仪征诗，有句「矮窗白纸出书声」174，反映的是农家子弟在读书，可见扬州人读书的普遍性。在安定、梅花、乐仪书院等教育机构配合下，扬州产生大量人才，更以家族的形式表现出来。阮元在〈扬州府志事志氏族表图说三门记〉文中缕述扬州府的人才家族:「江都、甘泉、仪征之耿氏、唐氏、杨氏、常氏、郑氏、秦氏、许氏、阮氏，兴化李氏，高邮王氏，宝应刘氏、乔氏，泰州宫氏等族，各以宰相九卿勋爵督抚等官，家自为谱。」175可见扬州文化渊源有自。

重视教育，是我国社会、家庭文化传统，极具传承价值。就此我亦想到两点，其一为提倡家庭教育，是德智体的全面教育，并非只是智力教育。其二是办好高等学校。盛清时代扬州教育基础最好，是官民齐努力，富商来赞助形成的，赞助者基本上是自觉自愿的，认识到自家的社会责任，徽商、西商深知必须回报土着，回报社会。这方面，马曰琯、程梦星、江春是榜样。有条件的人到扬州办学校，将会令人称羡。扬州若能成为大学城，岂不美哉！

扬州可以在文化教育、休闲文化方面再创辉煌，引领潮流。

\vfill
\begin{BookNotes}
\BookNote{172.}{《望江南百调》，第1043页。}
\BookNote{173.}{《汉书·董仲舒传》。}
\BookNote{174.}{《郑板桥全集·真州杂诗八首并及左右江县》，第127页。}
\BookNote{175.}{阮元文见《续纂扬州府志》卷23〈艺文下〉，第3册第1357页。}
\end{BookNotes}
